\section{Parameter Siklus Pembaruan (H3 dan H4)}

\begin{longtable}{
    @{}
    >{\raggedright\arraybackslash}p{5.0cm}          % Parameter
    >{\raggedleft\arraybackslash}p{3.0cm}           % Nilai
    >{\raggedright\arraybackslash}p{7.0cm}          % Keterangan
    @{}
}
\caption{Parameter simulasi siklus pembaruan lintas jendela} \\
\toprule
Parameter & Nilai & Keterangan \\
\midrule
\endfirsthead
\caption{Parameter simulasi siklus pembaruan lintas jendela (lanjutan)} \\
Parameter & Nilai & Keterangan \\
\midrule
\endhead
\bottomrule
\endfoot

Ambang pemicu drift ($\tau$) & $0,01$ & Struktur Chow-Liu dipelajari ulang bila divergensi Jensen-Shannon $> \tau$ \\
Jumlah jendela & 5 & Log dibagi menjadi lima jendela berurutan (empat siklus pembaruan) \\
Anggaran pembelajaran awal & 9.000 evaluasi & Jendela pertama / kondisi cold \\
Anggaran per pembaruan & 3.000 evaluasi & Tiap siklus pembaruan (warm/always) \\
Anggaran per aturan & 3.000 evaluasi & Batas per aturan pada siklus pembaruan \\
Ukuran populasi virtual ($NP$) & 100 & Sama seperti eksperimen H1 \\

\end{longtable}